\documentclass[10pt,a4paper]{article}
\usepackage[a4paper, margin=1.5cm]{geometry}
\usepackage[version=4]{mhchem}
\usepackage{chemfig}
\usepackage{xcolor}
\usepackage{amsmath}
\usepackage{amssymb}

\begin{document}

\section*{Fiche d'exercices -- séquence n°3~: Réactions d'oxydo-réduction}

\subsection*{EXERCICE 1~: QCM sur l'oxydo-réduction}

\begin{enumerate}
    \item Le nombre d'oxydation de l'élément cuivre dans l'ion cuivre (II) est~:
    \begin{itemize}
        \item[$\square$] 0
        \item[$\square$] -II
        \item[$\square$] +I
        \item[$\square$] +II
    \end{itemize}
    \textcolor{red}{Le nombre d'oxydation du cuivre dans \ce{Cu^2+} est +II.}
    \textcolor{red}{Le cuivre dans l'ion \ce{Cu^2+} a perdu deux électrons, ce qui lui confère un nombre d'oxydation de +II.}

    \item Sachant que le nombre d'oxydation de l'élément \ce{H} est de +I, quel est le nombre d'oxydation de l'élément oxygène \ce{O} dans la molécule d'eau oxygénée \ce{H2O2(aq)}~?
    \begin{itemize}
        \item[$\square$] 0
        \item[$\square$] -I
        \item[$\square$] -II
        \item[$\square$] +I
    \end{itemize}
    \textcolor{red}{Dans \ce{H2O2}, l'oxygène a un nombre d'oxydation de -I.}
    \textcolor{red}{Dans \ce{H2O2}, chaque atome d'oxygène a un nombre d'oxydation de -I, car l'hydrogène a un nombre d'oxydation de +I et la molécule est neutre.}

    \item Sachant que le nombre d'oxydation de l'élément \ce{H} est de +I et celui de l'élément oxygène de -II, quel est le nombre d'oxydation de l'élément carbone dans l'ion hydrogénocarbonate \ce{HCO3-(aq)}~?
    \begin{itemize}
        \item[$\square$] +III
        \item[$\square$] 0
        \item[$\square$] +IV
        \item[$\square$] -IV
    \end{itemize}
    \textcolor{red}{Dans \ce{HCO3-}, le carbone a un nombre d'oxydation de +IV.}
    \textcolor{red}{En utilisant les nombres d'oxydation de \ce{H} (+I) et \ce{O} (-II), on calcule que le carbone doit avoir un nombre d'oxydation de +IV pour équilibrer la charge de l'ion.}

    \item Une réaction avec une perte d'électrons est appelée une~:
    \begin{itemize}
        \item[$\square$] réduction~;
        \item[$\square$] oxydation~;
        \item[$\square$] corrosion~;
        \item[$\square$] oxydo-réduction.
    \end{itemize}
    \textcolor{red}{Une réaction avec une perte d'électrons est une oxydation.}
    \textcolor{red}{Une oxydation est définie comme une réaction où une espèce chimique perd des électrons.}

    \item Une réaction avec un transfert d'électrons est appelée une~:
    \begin{itemize}
        \item[$\square$] réduction~;
        \item[$\square$] oxydation~;
        \item[$\square$] corrosion~;
        \item[$\square$] oxydo-réduction.
    \end{itemize}
    \textcolor{red}{Une réaction avec un transfert d'électrons est une oxydo-réduction.}
    \textcolor{red}{Une réaction d'oxydo-réduction implique un transfert d'électrons entre un oxydant et un réducteur.}
\end{enumerate}

\subsection*{EXERCICE 2~: Couples redox}

Compléter le tableau ci-dessous et écrire le couple redox correspondant. On donne le nombre d'oxydation de l'élément \ce{H} est de +I.

\begin{center}
    \begin{tabular}{|c|c|c|c|}
        \hline
        Formule & n.o & Oxydant ou Réducteur & Couple redox \\
        \hline
        \ce{H+} & +I & Oxydant & \ce{H+/H2} \\
        \hline
        \ce{H2} & 0 & Réducteur & \ce{H+/H2} \\
        \hline
        \ce{I-} & -I & Réducteur & \ce{I2/I-} \\
        \hline
        \ce{I2} & 0 & Oxydant & \ce{I2/I-} \\
        \hline
        \ce{O2} & 0 & Oxydant & \ce{O2/H2O2} \\
        \hline
        \ce{H2O2} & -I & Oxydant/Réducteur & \ce{O2/H2O2} \\
        \hline
        \ce{Fe^2+} & +II & Réducteur & \ce{Fe^3+/Fe^2+} \\
        \hline
    \end{tabular}
\end{center}

\textcolor{red}{
\begin{itemize}
    \item \ce{H+/H2} avec n.o(\ce{H}) = +I et 0
    \item \ce{I2/I-} avec n.o(\ce{I}) = 0 et -I
    \item \ce{O2/H2O2} avec n.o(\ce{O}) = 0 et -I
    \item \ce{Fe^3+/Fe^2+} avec n.o(\ce{Fe}) = +III et +II
\end{itemize}}
\textcolor{red}{Les couples redox sont déterminés par les espèces oxydantes et réductrices. L'oxydant est l'espèce qui gagne des électrons, tandis que le réducteur est celle qui en perd.}

\subsection*{EXERCICE 3~: Demi-réactions}

Pour chacun des couples identifiés dans l'exercice 2, écrire l'équation de demi-réaction correspondante. NB~: Certaines réactions se font en milieu acide.

\begin{center}
    \begin{tabular}{|c|c|}
        \hline
        Couple redox & Demi-réaction \\
        \hline
        \ce{H+/H2} & \ce{2 H+ + 2 e- -> H2} \\
        \hline
        \ce{I2/I-} & \ce{I2 + 2 e- -> 2 I-} \\
        \hline
        \ce{O2/H2O2} & \ce{O2 + 2 H+ + 2 e- -> H2O2} \\
        \hline
        \ce{Fe^3+/Fe^2+} & \ce{Fe^3+ + e- -> Fe^2+} \\
        \hline
    \end{tabular}
\end{center}


\subsection*{EXERCICE 4~: Réaction entre les ions or (III) et le zinc}

Les ions or (III) réagissent avec le métal zinc pour donner un dépôt d'or métallique et des ions zinc (II). Les deux couples mis en jeu sont~: \ce{Au^3+/Au} et \ce{Zn^2+/Zn}.

\begin{enumerate}
    \item En déduire la position des deux couples sur une échelle de potentiel. \\
    \textcolor{red}{Le couple \ce{Au^3+/Au} a un potentiel standard plus élevé que \ce{Zn^2+/Zn}.}
    \textcolor{red}{L'or est un oxydant plus fort que le zinc, donc son potentiel standard est plus élevé.}

    \item Écrire les deux équations de demi-réaction correspondantes en identifiant l'oxydation et la réduction. \\
    \begin{itemize}
        \item Réduction~: \ce{Au^3+ + 3 e- -> Au}
        \item Oxydation~: \ce{Zn -> Zn^2+ + 2 e-}
    \end{itemize}

    \item Écrire alors l'équation de la réaction d'oxydo-réduction correspondante. \\
    \textcolor{red}{\ce{2 Au^3+ + 3 Zn -> 2 Au + 3 Zn^2+}}
    \textcolor{red}{L'équation est équilibrée en électrons et en masse.}
\end{enumerate}

\subsection*{EXERCICE 5~: Réaction entre le sulfate de cuivre (II) et le zinc}

On verse une solution de sulfate de cuivre (II) dans un erlenmeyer. On ajoute de la poudre de zinc et on maintient une agitation régulière pendant quelques instants. On filtre~: la solution limpide obtenue est incolore. La poudre ainsi recueillie est recouverte d'un dépôt métallique rouge.

\begin{enumerate}
    \item Schématiser l'expérience en utilisant le matériel et le vocabulaire adapté et en indiquant l'évolution de la couleur du milieu réactionnel. \\
    \textcolor{red}{Schéma à réaliser avec un erlenmeyer contenant une solution bleue de sulfate de cuivre (II) et de la poudre de zinc. Après réaction, la solution devient incolore et un dépôt rouge de cuivre apparaît.}

    \item Écrire l'équation de la réaction d'oxydo-réduction modélisant la transformation chimique lors de cette expérience. \\
    \textcolor{red}{\ce{Cu^2+ + Zn -> Cu + Zn^2+}}

    \item Identifier l'oxydant et le réducteur mis en jeu dans cette réaction. \\
    \textcolor{red}{L'oxydant est \ce{Cu^2+} et le réducteur est \ce{Zn}.}

    \item Déterminer l'entité chimique qui subit l'oxydation et celle qui subit la réduction. \\
    \textcolor{red}{Le zinc subit l'oxydation et le cuivre subit la réduction.}
\end{enumerate}
\end{document}
